\section{Supporting Lemmas}
\label{app:lemmas}

The proof core separates structural decoding from probabilistic
accounting. The well-formed transcript language of
\Cref{def:wf-transcripts} fixes which flattened queries are valid \TW{}
construction prefixes.
\Cref{lem:duplex-flat,lem:transcript-inj,lem:bridge-decode} then show
how such a query is flattened, parsed as a typed transcript, and
assigned an output point; the exact keyed decoder of
\Cref{def:keyed-twout} then decodes direct queries to candidate keys,
when any.
\Cref{lem:key-tests} turns those decoded keys into an adaptive key-test
bound. The remaining proof interfaces expose the events their consumers
must control: encryption-side freshness, verification hidden-key exposure
and root tag guessing, leaf tag collisions, and the committing
candidate-sequence collision event. Thus
the mu-ind and committing proofs only have to derive the relevant local
event premise before applying the corresponding accounting argument.

\paragraph{How to read this appendix.}
\Cref{fig:lemma-dag} records the proof dependency map. It groups local
interfaces when the exact proof citations would obscure the spine. The
Bridge \& Decode node is \Cref{lem:bridge-decode} together with the
exact keyed decoder of \Cref{def:keyed-twout}. The Root/Output
Separation node groups Output-Point Separation (\Cref{cor:output-sep}),
Opening Triple Separation (\Cref{cor:triple-sep}), and Final-Root Input
Separation (\Cref{prop:final-root-sep}); the use of Leaf-Transcript
Recovery (\Cref{lem:leaf-recovery}) is folded into the final-root
separation fact. The Key-Indexed Renaming node (\Cref{lem:key-renaming})
supplies the key-field disjointness and renaming used by hidden-key
encryption and verification accounting. The Adaptive Candidate node
(\Cref{lem:adaptive-candidate}) is the generic probabilistic engine: it
bounds the Leaf Tag Collision lemma (\Cref{lem:leaf-collision}) and,
through the Root Tag Candidate Sequence and Root Tag Hash and Target-Hit Bound
(\Cref{lem:root-candidate-seq,lem:root-tag-hash}) of \Cref{sec:cmtds3},
the root tag hash and committing proofs. The mu-ind proof reuses the
same structural lemmas through the explicit hybrid game sequence of
\Cref{def:adjacent-hybrid-games}, adding freshness, key-test,
leaf-collision, and verification-first root guessing accounting.

\begin{figure}[H]
  \centering
  \resizebox{\linewidth}{!}{%
  \begin{tikzpicture}[
    every node/.style={font=\footnotesize},
    lemma/.style={draw, rounded corners=2pt, align=center, inner sep=3pt,
                  minimum width=2.6cm, minimum height=0.7cm, fill=white},
    group/.style={draw, rounded corners=2pt, align=center, inner sep=3pt,
                  minimum width=2.8cm, minimum height=0.7cm, fill=white},
    consumer/.style={draw, dashed, rounded corners=2pt, align=center,
                     inner sep=3pt, minimum width=2.3cm, minimum height=0.7cm,
                     fill=white},
    arr/.style={-{Latex[length=1.6mm]}, semithick, draw=black!70},
  ]
    % Column 1: foundational lemmas.
    \node[lemma] (df) at (0,  2.6) {Duplex\\ Flattening};
    \node[lemma] (ba) at (0, -2.6) {Bridge \&\\ Decode};

    % Column 2: structural interfaces.
    \node[lemma] (ti) at (3.8,  2.6) {Transcript\\ Injectivity};
    \node[lemma] (kr) at (3.8,  0.8) {Key-Indexed\\ Renaming};
    \node[lemma] (kt) at (3.8, -1.0) {Adaptive\\ Key Tests};
    \node[group] (rs) at (3.8, -2.8) {Root/Output\\ Separation};

    % Column 3: local accounting lemmas.
    \node[lemma] (tf) at (7.6,  2.8) {Transcript\\ Freshness};
    \node[lemma] (lc) at (7.6,  1.2) {Leaf Tag\\ Collision};
    \node[lemma] (rg) at (7.6, -0.4) {Verification-First\\ Root Guessing};
    \node[lemma] (cc) at (7.6, -2.0) {Adaptive\\ Candidate};

    % Column 4: grouped proof spines.
    \node[group] (mi) at (11.4,  1.7) {Mu-ind\\ Accounting};
    \node[lemma] (rc) at (11.4, -0.5) {Root Tag\\ Candidate Seq.};
    \node[lemma] (rh) at (11.4, -2.0) {Root Tag\\ Hash/Target};

    % Column 5: consumers.
    \node[consumer] (c6) at (15.2,  1.7) {Mu-ind\\proof};
    \node[consumer] (c8) at (15.2, -1.3) {Hash +\\committing proofs};
    \node[consumer] (cb) at (15.2, -3.0) {App.~\ref{app:flat-sponge-lift}};

    \begin{scope}[on background layer]
      % Foundational -> structural interfaces.
      \draw[arr] (df) -- (ti);
      \draw[arr] (ti) -- (kr);
      \draw[arr] (ba.east) to[out=35, in=205] (kr.west);
      \draw[arr] (ba.east) to[out=10, in=200] (kt.west);
      \draw[arr] (ti.east) to[out=-35, in=145] (rs.west);
      \draw[arr] (ba) -- (rs);

      % Local accounting from the structural interfaces.
      \draw[arr] (ti.east) to[out=20, in=180] (tf.west);
      \draw[arr] (ti.east) to[out=0, in=180] (lc.west);
      \draw[arr] (ti.east) to[out=-20, in=170] (rg.west);
      \draw[arr] (cc) -- (lc);
      \draw[arr] (rs.east) to[out=25, in=190] (rc.west);
      \draw[arr] (ba.east) to[out=-5, in=200] (rc.west);
      \draw[arr] (rc) -- (rh);
      \draw[arr] (cc.east) to[out=-5, in=180] (rh.west);

      % Mu-ind accounting and consumer.
      \draw[arr] (tf.east) to[out=0, in=170] (mi.west);
      \draw[arr] (lc.east) to[out=0, in=185] (mi.west);
      \draw[arr] (rg.east) to[out=20, in=205] (mi.west);
      \draw[arr] (kr.east) to[out=25, in=195] (mi.west);
      \draw[arr] (kt.east) to[out=35, in=215] (mi.west);
      \draw[arr] (mi) -- (c6);

      % Root tag hash and committing consumers.
      \draw[arr] (lc.east) to[out=-20, in=160] (c8.west);
      \draw[arr] (rs.east) to[out=-5, in=185] (c8.west);
      \draw[arr] (rh.east) -- (c8.west);

      % The flat sponge lift consumes bridge decoding directly.
      \draw[arr] (ba.east) to[out=-20, in=180] (cb.west);
    \end{scope}
  \end{tikzpicture}
  }
  \caption{Proof dependency map for the supporting components in this
    appendix and their consumers. Solid boxes are lemmas or grouped local
    proof interfaces; dashed boxes are proof components that consume them.}
  \label{fig:lemma-dag}
\end{figure}

\subsection{Duplex Flattening}
\label{app:lem-duplex-flat}

\begin{lemma}[Duplex Flattening]
\label{lem:duplex-flat}
Every \TW{} root or leaf duplex output is equal to a \cite{BDPV11} sponge
output on the corresponding flattened duplex-call transcript, and the
ciphertext-absorbing message phase admits a flattened duplex transcript
that is well defined for both encryption and decryption.
\end{lemma}

\begin{proof}
Let $\mathsf{D}$ be a duplex object using permutation $\pi$, rate $r$,
and padding $\padtenone$. For a sequence of duplex calls $Z_i =
\mathsf{D}.\duplexing(\sigma_i, \ell_i)$ with each $\sigma_i$ short
enough to fit in one duplex block after padding, \cite[Lemma~3]{BDPV11}
gives
\[
  Z_i = \Sponge[\pi, \padtenone, r]
    \bigl(\sigma_0 \concat \mathsf{pad}_0 \concat \sigma_1
          \concat \mathsf{pad}_1 \concat \cdots \concat \sigma_i,\, \ell_i\bigr),
\]
where $\mathsf{pad}_j = \padtenone[r](|\sigma_j|)$, and Lemma~4 says
that, for fixed $\padtenone$ and $r$, the map from the duplex input-string
sequence to the flattened sponge input is injective.

In \TW{}, each duplex call absorbs a byte-aligned block $\sigma$
followed by a $4$-bit phase suffix $\sigph \in \{\siginit,
\sigadmore, \dots, \sigagglast\}$
(\Cref{sec:domain-sep}), so the duplex-call input is
$\sigma \concat \sigph$ followed by $\padtenone$. With
$r = 1344$ and $\rho_{\max}(\padtenone, r) = r - 2 = 1342$ bits, the
maximum byte-aligned input is $\rho = 167$ bytes
(\Cref{sec:params}), giving a worst-case duplex-call input of $8 \cdot
167 + 4 = 1340 \leq 1342$ bits. The two initialization calls are
shorter still: $|\prfxroot \concat K \concat U| + 4 = 8
\cdot (6 + 32 + 32) + 4 = 564$ bits and $|\prfxleaf \concat
K \concat U \concat \langle j \rangle_8| + 4 = 8 \cdot (6 + 32 + 32 +
8) + 4 = 628$ bits. \cite[Lemmas~3 and~4]{BDPV11} therefore apply
independently to every \TW{} duplex object.

The root duplex is initialized by absorbing $\prfxroot
\concat K \concat U$ under \siginit (\Cref{alg:enc}), then
absorbs the associated data blocks under \sigadmore /
\sigadlast, the root ciphertext blocks in the message phase under
\sigmsgmore / \sigmsglast, and the leaf tag aggregation
blocks under \sigaggmore / \sigagglast. By the
duplex-to-sponge equality, every root duplex output is the sponge
function evaluated on the root transcript prefix accumulated up to the
relevant output point. The same holds for each leaf duplex,
initialized by $\prfxleaf \concat K \concat U \concat
\langle j \rangle_8$ and absorbing only its own ciphertext blocks in the
message phase. In decryption and verification, \Cref{alg:dec} absorbs the
submitted ciphertext blocks before the root tag comparison, so the same
flattened message phase transcript is defined whether the comparison
accepts or rejects.
\end{proof}

\subsection{Well-formed Transcript Prefixes}
\label{app:wf-transcripts}

For a duplex-call prefix $((\sigma_0, \sigph[0]), \dots, (\sigma_i,
\sigph[i]))$, write
\[
  \Flat\bigl((\sigma_0, \sigph[0]), \dots, (\sigma_i, \sigph[i])\bigr)
  = \bigconcat_{j=0}^{i-1}\!
    \bigl(\sigma_j \concat \sigph[j] \concat \mathsf{pad}_j\bigr)
    \concat \sigma_i \concat \sigph[i],
\]
where $\mathsf{pad}_j = \padtenone[r](|\sigma_j \concat \sigph[j]|)$,
the padding of the $j$-th full duplex-call input.
The current call's padding $\mathsf{pad}_i$ is not part of $\Flat$;
it is applied internally by the sponge function. Equivalently, after
the $i$-th call the duplex state is the sponge state after absorbing
$\Flat(\cdot) \concat \mathsf{pad}_i$.

\begin{definition}[Well-Formed \TW{} Transcript Prefixes]
\label{def:wf-transcripts}
A typed \TW{} transcript prefix is \emph{well formed} if it is a prefix,
ending at a construction transcript lookup point, of one of the following
two call languages.
\begin{itemize}
\item A root transcript starts with
  $(\prfxroot \concat K \concat U,\siginit)$, with fixed-width fields
  $K \in \bits^k$ and $U \in \bits^\nu$. It then absorbs the
  associated data byte string (only when the associated data is
  nonempty; the phase is otherwise elided, \Cref{sec:enc}), the root
  ciphertext byte string, and the aggregation byte string (only when
  $n \geq 1$; when $n = 0$ the aggregation phase is elided,
  \Cref{sec:enc}, and the root tag is emitted by the final
  $\sigmsglast$ message call), in that order. Each present phase byte
  string is partitioned by the block rule
  of \Cref{alg:helpers}: the empty string gives one empty final block,
  and a nonempty string gives zero or more $\rho$-byte non-final blocks
  followed by one final block. The phases use respectively
  $\sigadmore/\sigadlast$, $\sigmsgmore/\sigmsglast$, and
  $\sigaggmore/\sigagglast$.
\item A leaf-$j$ transcript starts with
  $(\prfxleaf \concat K \concat U \concat \langle j \rangle_8,\siginit)$,
  where $1 \leq j < 2^{64}$ and the fields have fixed width. It then
  absorbs the leaf ciphertext byte string under
  $\sigmsgmore/\sigmsglast$, using the same block rule.
\end{itemize}
The message byte string in this language is always the ciphertext byte
string absorbed by the construction: encryption obtains it by XORing the
plaintext with the current keystream, while decryption and verification
use the submitted ciphertext. Thus a verification transcript is well
formed for every in-domain ciphertext before the tag comparison is made;
acceptance is not part of the transcript language.

The root aggregation byte string has the form
$T_1 \concat \cdots \concat T_n \concat \langle n \rangle_8$, where each
$T_j$ has fixed byte length $\tauleaf/8$. The trailing fixed-width field
therefore gives a unique parse of the leaf tag sequence. The
output-point class of a well-formed counted prefix is the class in
\Cref{tab:keyed-twout} determined by its duplex instance and final-call
suffix.
\end{definition}

\begin{remark}[Well-Formedness versus Realizability]
\label{rem:wf-vs-realizable}
The language of \Cref{def:wf-transcripts} is deliberately larger than
the set of transcript prefixes that \TW{} computations realize. For
example, it does not require an aggregation phase to follow a root
message phase of exactly $\beta$ bytes, although every real computation
with $n \geq 1$ absorbs a full $\beta$-byte root chunk (\Cref{alg:enc}).
The inclusion runs in the safe direction everywhere: every
construction-generated prefix is well formed, so the structural lemmas
of this appendix and the separation propositions of
\Cref{sec:root-tag-hash-security} apply to construction transcripts as
restrictions of statements proved on the larger language, while the
direct-query decoder of \Cref{def:keyed-twout} and the root tag
candidate predicate of \Cref{lem:root-candidate-seq} accept the larger
language and therefore only over-count adversary queries, which loosens
no bound.
\end{remark}

Name the tag-emitting transcript families:
\begin{align*}
  \Rtag   &= \text{root transcript prefix whose sponge output gives the final root tag}, \\
  \Ltag(j)   &= \text{leaf-$j$ transcript prefix whose sponge output gives the leaf tag}.
\end{align*}
The final root tag is emitted by the $\sigagglast$ aggregation call (the
$\Rtag$ family) when $n \geq 1$, and by the closing $\sigmsglast$
message call (an $\Rmsg^+$ prefix, in the notation of
\Cref{cor:output-sep}) when $n = 0$ and the aggregation phase is elided.
Since $n = \leaves(C)$ is fixed by $\|C\|$, every computation on a given
ciphertext emits its tag at the output point selected by $\|C\|$.
Intermediate construction transcript lookup points (init, non-final
AD, non-final message, and non-final aggregation) are also typed
prefixes covered by \Cref{lem:bridge-decode}; the exact keyed
direct-query decoder names them in \Cref{tab:keyed-twout}.

\begin{definition}[Root Tag Output Point and Final Root Transcript Class]
\label{def:root-tag-output-point}
For any in-domain encryption or verification computation on ciphertext
$C$, the \emph{root tag output point} is the construction lookup that
emits the root tag: $\Rtag$ when $\leaves(C) \geq 1$, and the closing
root message point $\Rmsg^+$ when $\leaves(C) = 0$ and aggregation is
elided. The \emph{final root transcript class} of such a computation is
the set of construction computations that realize the same bridged final
root input at this root tag output point. A class is
\emph{encryption-first} if its first construction generation is by an
encryption query, and \emph{verification-first} if its first construction
generation is by a non-replay verification query.
\end{definition}

\subsection{Transcript Injectivity}
\label{app:lem-transcript-inj}

\begin{lemma}[Transcript Injectivity]
\label{lem:transcript-inj}
Well-formed \TW{} flattened duplex-call transcript prefixes are
injectively encoded and separated by duplex instance (root vs leaf),
leaf index $j$, phase suffix, and duplex-call order. For final-root
prefixes, equality of flattened inputs recovers the same root
initialization, AD blocks, root ciphertext blocks, aggregation input,
and leaf tag sequence.
\end{lemma}

\begin{proof}
By \Cref{lem:duplex-flat}, every well-formed output point can be written
as a \cite{BDPV11} flattened sponge input for the corresponding
duplex-call prefix, namely $\Flat$ as defined above. The deterministic
current-call $\padtenone$ is applied internally by the sponge function
and is not part of the input to which \cite[Lemma~4]{BDPV11} is applied.
That lemma recovers the duplex-call sequence from the flattened sponge
input; in \TW{}, each call input is
$\sigma \concat \sigph$ with byte-aligned $\sigma$ and
fixed $4$-bit suffix $\sigph$, so recovering the call input recovers the pair
$(\sigma, \sigph)$.

\Cref{def:wf-transcripts} fixes the root and leaf initialization fields,
the admissible phase order, and the final-call suffixes. The distinct
prefixes $\prfxroot$ and $\prfxleaf$ separate root from leaf transcripts;
the fixed-width leaf field $\langle j \rangle_8$ separates leaves; and
the seven suffix values of \Cref{tab:suffixes} separate phases and
non-final/final calls. Thus equality of flattened inputs forces equality
of the recovered instance type, key, nonce, leaf index when present,
phase suffixes, and duplex-call order.

\paragraph{Phase-recovery sub-claim.}
For each present phase $P \in \{\mathsf{ad}, \mathsf{msg}, \mathsf{agg}\}$
(the $\mathsf{ad}$ phase is present exactly when the associated data is
nonempty; see \Cref{sec:enc}),
\TW{} absorbs the phase's input byte string $s_P$ by partitioning into
$\rho$-byte blocks (with the empty case yielding a single empty block,
per \Cref{sec:enc}) and feeding each block to the duplex with suffix
$\sigma_P^-$ (non-final) or $\sigma_P^+$ (final).
\Cref{lem:duplex-flat} and \cite[Lemma~4]{BDPV11} recover the
corresponding $(\sigma, \sigph)$ sequence from the flattened input.
Because the partition map is deterministic and uniquely invertible on
its image (the empty case yielding the single empty block carrying
$\sigma_P^+$), concatenating the recovered $\sigma$ blocks
reconstructs $s_P$ exactly. The map from $s_P$ to the phase's $(\sigma,
\sigph)$ sequence is therefore injective.

Applied phase by phase to a well-formed final-root transcript:
\begin{itemize}
\item AD. The associated data phase contributes a $\sigadlast$-tagged
  block exactly when the associated data is nonempty; an empty
  associated data elides the phase and contributes no block. When the
  phase is present the sub-claim recovers the AD byte string. Thus
  distinct associated data strings $A_i \neq A_j$ are separated either
  by the presence versus absence of the AD phase (when exactly one is
  empty) or, when both are nonempty, by their distinct AD
  $\sigma$-block sequences.
\item MSG. Distinct absorbed root or leaf ciphertext byte strings
  produce distinct MSG $\sigma$-block sequences; the message language of
  \Cref{def:wf-transcripts} is ciphertext-absorbing for encryption and
  decryption alike.
\item AGG. The aggregation phase is present exactly when $n \geq 1$
  (\Cref{sec:enc}). When it is present, the sub-claim applied to the
  aggregation byte string $T_1 \concat \cdots \concat T_n \concat
  \langle n \rangle_8$ recovers that byte string. The final eight bytes
  give $\langle n \rangle_8$, and each leaf tag has fixed byte length
  $\tauleaf / 8$, so the byte string has a unique parse as a leaf tag
  sequence. Two root aggregation encodings are therefore equal exactly
  when their leaf tag sequences are equal. When
  $n = 0$ the aggregation phase is absent and the final-root transcript
  ends at the recovered $\sigmsglast$ root message call; the MSG bullet
  recovers the root ciphertext byte string, and the absence of any
  $\sigaggmore/\sigagglast$ call distinguishes such a prefix from every
  $n \geq 1$ final-root prefix.
\end{itemize}

Therefore the recovered final-root transcript determines the root
initialization, AD blocks, root ciphertext blocks, and, when
$n \geq 1$, the aggregation byte string and leaf tag sequence.
\end{proof}

\begin{lemma}[Leaf-Transcript Recovery]
\label{lem:leaf-recovery}
Let two construction-generated final-root prefixes have equal flattened
inputs, and suppose that, for each position $j$ of their common leaf tag
sequence, the two computations' position-$j$ final leaf transcripts are
either equal or receive distinct leaf tags. Then equality of flattened
inputs recovers the same leaf transcript sequence, and therefore the
same full absorbed ciphertext. Both computations' position-$j$ final
leaf transcripts are leaf-$j$ transcripts under the common recovered
$(K, U)$, so the hypothesis holds whenever no two distinct
construction-generated final leaf transcripts with the same key, nonce,
and leaf index receive the same leaf tag; when one of the two prefixes
is generated by an encryption computation, it suffices that no
construction-generated final leaf transcript receives the same leaf tag
as a distinct encryption-generated final leaf transcript with the same
key, nonce, and leaf index.
\end{lemma}

\begin{proof}
By Transcript Injectivity (\Cref{lem:transcript-inj}), the two equal
flattened final-root inputs recover the same root initialization, AD
blocks, root ciphertext blocks, and, when $n \geq 1$, the same
aggregation byte string and leaf tag sequence; the trailing
$\langle n \rangle_8$ field fixes the common leaf count $n$. When
$n = 0$ no leaf is present and the recovered root ciphertext blocks
already constitute the full absorbed ciphertext. When $n \geq 1$, the
aggregation byte string parses uniquely as the leaf tag sequence
$T_1 \concat \cdots \concat T_n$; the two leaf tag sequences are equal,
so for each $j$ the two computations' position-$j$ final leaf
transcripts---leaf-$j$ transcripts carrying the common recovered
$(K, U)$ (\Cref{def:wf-transcripts}, \Cref{sec:enc})---receive the same
leaf tag, and by hypothesis are equal, hence so are
their absorbed leaf ciphertext bodies. Together with the common root
ciphertext blocks this recovers the same full absorbed ciphertext.
\end{proof}

\subsection{Bridge and Typed Decoding}
\label{app:lem-bridge-decode}

\begin{lemma}[Bridge Injectivity]
\label{lem:bridge-decode}
The random oracle games used in
\Cref{sec:aead-security,sec:root-tag-hash-security} and \Cref{cor:committing-combined} are
ordinary random oracle games on a single random oracle
$\ROflat : \bits^* \to \bits^\infty$ over the unpadded
\cite{BDPV08} $H$-input domain (\Cref{sec:ro-model}), with \TW{} root and
leaf duplex outputs answered by the typed restriction
\[
  \RF(X) = \ROflat(\flatmap(X)),
  \qquad \flatmap(X) = \bridge[r, r_{\max}]\bigl(\mathsf{raw\_flat}(X)\bigr),
\]
for $r = r_{\max} = 1344$. Here $\mathsf{raw\_flat}(X)$ is the native
flattened input $\Flat(\cdot)$ of \Cref{app:wf-transcripts} applied to
the duplex-call sequence of $X$, and $\bridge[r, r_{\max}]$ is the
\cite{BDPV11} Theorem~3 preprocessing map. The induced map
$X \mapsto \flatmap(X)$ is injective on well-formed typed \TW{} duplex
transcript prefixes.
\end{lemma}

\begin{proof}
The bridge $\bridge[r, r_{\max}]$ is described by the proof of
\cite[Theorem~3]{BDPV11} in three steps: (1)~pad the input with
$\padtenone[r]$; (2)~for each $r$-bit block, append $r_{\max} - r$
zero bits; (3)~strip the trailing $\padten$ from the result. The
$\padtenone$ padding in step~(1) appends a~$1$, then
$q = (-|\mathsf{raw\_flat}(X)| - 2) \bmod r$ zero bits, then a closing
$1$; the $\padten$-unpadding in step~(3) strips the $r_{\max} - r$
trailing zero bits introduced by step~(2) together with the closing
$1$ of step~(1). For \TW{} with $r = r_{\max}$, step~(2) appends an
empty string and step~(3) strips zero trailing zeros and the closing
$1$, so the net effect of steps~(1) and~(3) is to map a native flattened
input $W = \mathsf{raw\_flat}(X)$ to $W \concat 1 \concat 0^q$. Any
string in this image has a unique inverse: $q$ is the number of trailing
zero bits, the preceding bit is the inserted $1$, and deleting this
suffix recovers $W$.

Injectivity of $\flatmap$ on well-formed typed \TW{} prefixes then
follows from \Cref{lem:transcript-inj}: distinct well-formed typed
prefixes have distinct
$\mathsf{raw\_flat}(\cdot)$ images, and the bridge is injective on
that image. Direct adversary queries on inputs outside the bridged
\TW{} well-formed transcript language still count in $\qRO(\mathcal{B})$
(\Cref{sec:resources}) but cannot trigger
$\badkey$.
\end{proof}

The bridge injectivity established here underlies Opening Triple
Separation (\Cref{cor:triple-sep}, \Cref{sec:cmtds3}): distinct
$(K, U, A)$ triples yield distinct bridged final root transcripts.

\begin{corollary}[Output-Point Separation]
\label{cor:output-sep}
Every well-formed \TW{} typed transcript prefix at a counted
construction transcript lookup point belongs to exactly one of the
following output-point classes, identified by the recovered duplex
instance and the $4$-bit suffix of its final call:
\begin{itemize}
\item \Rinit (root duplex, suffix \siginit),
\item \Linit($j$) (leaf-$j$ duplex, suffix \siginit),
\item $\Rad^-$ / $\Rad^+$ (root duplex, suffix \sigadmore or \sigadlast),
\item $\Rmsg^-$ / $\Rmsg^+$ (root duplex, suffix \sigmsgmore or \sigmsglast),
\item $\Lmsg^-(j)$ (leaf-$j$ duplex, suffix \sigmsgmore),
\item $\Ltag(j)$ (leaf-$j$ duplex, suffix \sigmsglast),
\item \Ragg (suffix \sigaggmore),
\item $\Rtag$ (root duplex, suffix \sigagglast).
\end{itemize}
For any two such prefixes whose bridged flattened inputs are equal, the
separation of \Cref{lem:transcript-inj} recovers the same duplex
instance, leaf index $j$ when present, $(K, U)$ pair, output-point
class, and duplex-call sequence. Equivalently, any two distinct counted
transcript prefixes already differ in one of these recovered fields, and
hence in their flattened inputs.
\end{corollary}

\begin{proof}
\Cref{lem:bridge-decode} recovers equality of native flattened inputs
from equality of bridged inputs, after which \Cref{lem:transcript-inj}
recovers the duplex-call sequence and its per-call
$(\sigma, \sigph)$ pairs. The first call's $\sigma$ is
$\prfxroot \concat K \concat U$ for root instances and
$\prfxleaf \concat K \concat U \concat \langle j \rangle_8$ for leaf
instances, with fixed-width fields, so it fixes the instance type (by
$\prfxroot$ vs $\prfxleaf$), the $(K, U)$ pair, and the leaf index $j$;
within each instance type the phase suffixes are pairwise distinct, so the
recovered instance type together with the final-call suffix fixes the
output-point class, and each prefix therefore lies in exactly one class.
Conversely, two distinct well-formed prefixes in the same class on the
same instance differ in some recovered call's $(\sigma, \sigph)$ pair
or in the
sequence length, hence in their flattened inputs.
\end{proof}

\begin{definition}[Exact Keyed Decoder]
\label{def:keyed-twout}
For a direct query to $\ROflat$ with input $Y$ and finite length
$\ell$, define $\KeyedTWOut(Y)$ as a partial map from $\bits^*$ to
$\bits^k \times \mathcal{C}_{\mathsf{out}}$. The map ignores
$\ell$. It returns $(K,\mathcal{C})$ iff
there exists a typed \TW{} duplex-call prefix $X$ such that
$Y = \flatmap(X)$, the raw flattened prefix of $X$ parses as a
well-formed \TW{} transcript prefix in the class $\mathcal{C}$ listed in
\Cref{tab:keyed-twout}, and the first call of $X$ contains the
candidate key $K$ in the fixed-width root field
$\prfxroot \concat K \concat U$ or leaf field
$\prfxleaf \concat K \concat U \concat \langle j \rangle_8$. If no
such prefix exists, $\KeyedTWOut(Y) = \bot$. In particular, inputs with
the right final suffix but a malformed phase order, noncanonical block
partition, malformed aggregation string, or out-of-range leaf index are
outside the well-formed language and decode to $\bot$. A non-$\bot$
decoding always fixes the candidate key $K$ and nonce $U$ in the
fixed-width initialization field. It fixes the associated data $A$ only
for root prefixes whose recovered phase order determines the complete AD
string: a final AD prefix, or any later root message or aggregation
prefix, with absence of AD calls determining $A=\varepsilon$. We use the
key projection $(K,\mathcal{C})$ throughout. When a game hides a single
context component $Z \in \{K,U,A\}$, write $\KeyedTWOut_Z(Y)$ for the
decoded value of that component when it is determined, and $\bot$
otherwise.
\end{definition}

By \Cref{lem:bridge-decode,cor:output-sep}, $\KeyedTWOut$ is well
defined: a bridged input cannot parse as two distinct counted
well-formed prefixes, and an accepted input fixes a unique candidate key
and output-point class. Inputs outside this well-formed transcript
language decode to $\bot$.

\begin{table}[ht]
  \centering
  \small
  \begin{tabular}{l l l}
    \toprule
    Class $\mathcal{C}$ & Instance & Final call / role \\
    \midrule
    \Rinit       & root     & \siginit, root initialization \\
    $\Rad^-$     & root     & \sigadmore, non-final AD block \\
    $\Rad^+$     & root     & \sigadlast, final AD block \\
    $\Rmsg^-$    & root     & \sigmsgmore, non-final root message block \\
    $\Rmsg^+$    & root     & \sigmsglast, final root message block \\
    \Ragg        & root     & \sigaggmore, non-final aggregation block \\
    $\Rtag$      & root     & \sigagglast, final aggregation / root tag \\
    \Linit($j$)  & leaf $j$ & \siginit, leaf initialization \\
    $\Lmsg^-(j)$ & leaf $j$ & \sigmsgmore, non-final leaf message block \\
    $\Ltag(j)$   & leaf $j$ & \sigmsglast, final leaf message block / leaf tag \\
    \bottomrule
  \end{tabular}
  \caption{Counted keyed transcript classes for $\KeyedTWOut$.}
  \label{tab:keyed-twout}
\end{table}

The table is a transcript-lookup table, not an observation-length
table. A class is counted even when the construction's requested
output length at that call is zero, including root initialization,
non-final AD and aggregation calls, a final AD call whose following
root chunk requests no keystream bytes, and the final root message
call. A direct query triggers the same decoder regardless of $\ell$:
short finite-output queries and zero-length queries $(Y,0)$ still
count as direct queries to the address $Y$. Leaves have no AD or
aggregation classes; a leaf transcript enters the table only through
leaf initialization, non-final leaf message blocks, or the final leaf
tag point. When the associated data is empty the root transcript
likewise has no $\Rad^-$ or $\Rad^+$ class (\Cref{sec:enc}); the
\Rinit class then supplies the first root keystream block, exactly as
\Linit($j$) supplies the first leaf keystream block, and the decoder
still ignores $\ell$. Symmetrically, when $n = 0$ the root transcript
has no \Ragg or $\Rtag$ class (\Cref{sec:enc}); the final root message
call $\Rmsg^+$ then supplies the root tag, exactly as $\Ltag(j)$
supplies a leaf tag, and the decoder again ignores $\ell$. An $n \geq 1$
computation also reaches an $\Rmsg^+$ point when closing chunk $0$, but
there the construction requests zero output and continues into the
aggregation phase, so the $\Rmsg^+$ value is exposed as a tag only on
the $n = 0$ path; since $n = \leaves(C)$ is fixed by $\|C\|$, the two
uses never coincide on a single ciphertext. We call \Rinit, $\Rad^+$,
$\Rmsg^-$, \Linit($j$), and $\Lmsg^-(j)$ the \emph{stream output
points}: by the schedule just described, these are the only classes at
which a construction computation requests keystream bytes, while tag
values are exposed only at $\Ltag(j)$ and $\Rtag$ points and, on the
$n = 0$ path, at $\Rmsg^+$.

The hidden-key encryption, verification, and mu-ind accounting all
reindex transcript addresses by the key field. We record that reindexing
once.

\begin{lemma}[Key-Indexed Transcript Renaming]
\label{lem:key-renaming}
For a key $K \in \bits^k$, let $\mathcal{Y}_K$ be the set of bridged
inputs $\flatmap(X)$ of well-formed \TW{} transcript prefixes $X$ whose
first-call key field is $K$---the counted hidden-key transcript addresses
of \Cref{tab:keyed-twout} for key $K$. Then:
\begin{enumerate}
\item \emph{Disjointness.} For distinct keys $K_0 \neq K_1$, the sets
  $\mathcal{Y}_{K_0}$ and $\mathcal{Y}_{K_1}$ are disjoint.
\item \emph{Renaming.} The map $\Phi_{K_0 \to K_1}$ that replaces the
  first-call key field $K_0$ by $K_1$ and leaves every other recovered
  field unchanged---nonce, leaf index, phase sequence, AD blocks,
  ciphertext blocks, aggregated leaf tag bytes, and root/leaf
  instance---induces through $\flatmap$ a bijection
  $\flatmap \circ \Phi_{K_0 \to K_1} \circ \flatmap^{-1} :
  \mathcal{Y}_{K_0} \to \mathcal{Y}_{K_1}$ preserving output-point
  classes and equality relations among addresses.
\item \emph{Direct-query avoidance.} A direct query input
  $Y \in \mathcal{Y}_K$ has $\KeyedTWOut(Y) = (K,\mathcal{C})$ for a
  counted class $\mathcal{C}$; hence if no prior direct query decodes to
  $K$ then $\mathcal{Y}_K$ is disjoint from all prior direct-query
  inputs.
\end{enumerate}
\end{lemma}

\begin{proof}
By bridge injectivity (\Cref{lem:bridge-decode}) and Transcript
Injectivity (\Cref{lem:transcript-inj}), a bridged input $\flatmap(X)$
determines the native flattened prefix and hence the first duplex call
$\prfxroot \concat K \concat U$ or
$\prfxleaf \concat K \concat U \concat \langle j \rangle_8$, in
particular its fixed-width key field $K$. \emph{(1)} If
$\flatmap(X) = \flatmap(X')$ with key fields $K_0$ and $K_1$ then
$K_0 = K_1$, so distinct keys give disjoint address sets. \emph{(2)}
$\Phi_{K_0 \to K_1}$ rewrites only the key field, leaving the recovered
duplex-call sequence otherwise intact, so it is a bijection on
well-formed prefixes preserving the output-point class
(\Cref{cor:output-sep}) and the equality pattern among prefixes;
$\flatmap$ transports it to the stated bijection on addresses. \emph{(3)}
By \Cref{lem:bridge-decode} a $Y = \flatmap(X) \in \mathcal{Y}_K$ decodes
under $\KeyedTWOut$ to $(K,\mathcal{C})$; the contrapositive gives the
avoidance claim.
\end{proof}

\subsection{Adaptive Candidate Bounds}
\label{app:lem-candidate-collisions}

The Adaptive Candidate Collision and Preimage lemma is stated as
\Cref{lem:adaptive-candidate} in \Cref{sec:root-tag-hash-security},
where a proof sketch also appears; this subsection supplies the full
proof.

\begin{proof}[Proof of \Cref{lem:adaptive-candidate}]
Both parts expose the interaction in chronological order and sample each
candidate's $\ROflat$ value lazily at its append step, which the premises
license: when $Y_j$ is appended its string is fixed by the prior state
(predictability) and no bit of $\ROflat(Y_j)$ has yet been
revealed (unrevealed value), so $\ROflat(Y_j)$ may be drawn then, uniform
and independent of everything so far.

\emph{Collision.} For a fixed position set $I = \{i_1 < \cdots < i_s\}
\subseteq \{1,\dots,L\}$, let $E_I$ be the event that these candidate
positions exist and their $\tau$-bit projections are all equal. Condition
on any positive-probability prefix immediately before position $i_r$ is
appended, for $r \geq 2$, and on the values sampled at the earlier
positions in $I$. By the above, $\ROflat(Y_{i_r})$ is then uniform and
independent of the conditioning, so its $\tau$-bit projection equals the
already fixed projection of $Y_{i_1}$ with probability $2^{-\tau}$.
Multiplying over $r = 2,\dots,s$ gives $\Pr[E_I] \leq 2^{-\tau(s-1)}$; if
some position in $I$ is never appended then $E_I$ is false and the bound
still holds. A union bound over the $\binom{L}{s}$ position sets gives
the claim.

\emph{Preimage.} For a fixed position $j \in \{1,\dots,L\}$, let $E_j$
be the event that position $j$ exists and
$\bigl\lfloor \ROflat(Y_j) \bigr\rfloor_\tau = t$. Condition on any
positive-probability prefix immediately before $Y_j$ is appended; the
target $t$ is a constant. By the above, $\ROflat(Y_j)$ is uniform and
independent of the conditioning, so its $\tau$-bit projection equals $t$
with probability $2^{-\tau}$. If position $j$ is never appended, $E_j$ is
false. A union bound over the $L$ positions gives the claim.

\emph{Second-preimage corollary.} If the designated position $j_0$ is
not appended, the event is false. Otherwise split the other candidates
according to whether they occur before or after $j_0$. For earlier
candidates, condition on any positive-probability prefix immediately
before $Y_{j_0}$ is appended, including the already sampled projections
of candidates $1,\dots,j_0-1$. The value $\ROflat(Y_{j_0})$ is then
fresh and uniform, so its $\tau$-bit projection hits one of those at most
$j_0-1$ earlier projections with probability at most
$(j_0-1)2^{-\tau}$. For a later candidate $Y_j$, $j>j_0$, condition on
the full prefix immediately before it is appended, including the fixed
projection $t_0 = \lfloor\ROflat(Y_{j_0})\rfloor_\tau$. The value
$\ROflat(Y_j)$ is fresh and uniform at its append step, so it hits
$t_0$ with probability $2^{-\tau}$. A union bound over the at most
$L-j_0$ later candidates gives an additional $(L-j_0)2^{-\tau}$.
Adding the two sides gives
$\Pr[\mathsf{sec}_{\tau,j_0}] \leq (L-1) \cdot 2^{-\tau}$; the
designated candidate is excluded, so its trivially equal projection does
not enter the count.
\end{proof}

\subsection{Adaptive Key Tests}
\label{app:lem-key-tests}

\begin{lemma}[Adaptive Key-Test Bound]
\label{lem:key-tests}
Consider a \TW{} random oracle model game with a single hidden context
component $Z$ sampled uniformly from $\bits^w$ for some $w \geq k$---the
standard case being the hidden key, $Z = K$ and $w = k$, used in the
AE-style proofs---direct adversary access to
$\ROflat$, and an immediate abort event $\badkey$ that fires on a
direct query $(Y,\ell)$ exactly when $Y$ decodes to some counted
output-point class $\mathcal{C}$ of \Cref{tab:keyed-twout} with
$\KeyedTWOut_Z(Y) = Z$. Suppose that, before the
abort, the following invariant holds for every direct query position
$i$: after conditioning on the full visible execution prefix and on
the set $S_i$ of distinct decoded candidate $Z$-values from earlier direct
queries, the hidden $Z$ is uniform on $\bits^w \setminus S_i$. Then
\[
  \Pr[\badkey]
  \leq \qRO / 2^w
  \leq \KeyGuess_k(\qRO)
  = \qRO / 2^k,
\]
where $\qRO$ is the execution-uniform count of direct $\ROflat$ queries.
\end{lemma}

\begin{proof}
If $\qRO \geq 2^w$ the right-hand side $\qRO/2^w$ is at least one, so the
claim is trivial. Assume $\qRO < 2^w$.

Order the direct $\ROflat$ queries. By the definition of
$\KeyedTWOut_Z$ and \Cref{lem:bridge-decode}, each query input either
decodes to $\bot$ or decodes to a unique candidate $Z$-value and counted
output-point class. Queries that decode to $\bot$, and repeated
queries whose decoded $Z$-value already lies in the current set $S_i$, do
not change the set of possible first successful tests. Removing
them can only shorten the sequence, so it suffices to analyze the
subsequence of distinct fresh decoded candidate $Z$-values
$Z'_1, \dots, Z'_m$, where $m \leq \qRO$.

Let $F_i$ be the event that the first $i$ fresh decoded candidate values
all miss the hidden $Z$. The invariant says that, conditioned on the
visible prefix before the $(i+1)$-st fresh test and on $F_i$, the
hidden $Z$ is uniform over the $2^w - i$ values not yet tested. The next
fresh candidate value $Z'_{i+1}$ is one of those remaining values and is
determined by the adversary's next direct query, so
\[
  \Pr[Z = Z'_{i+1} \mid F_i] = \frac{1}{2^w - i},
  \qquad
  \Pr[F_{i+1} \mid F_i] = 1 - \frac{1}{2^w - i}.
\]
Multiplying the miss probabilities gives
\[
  \Pr[F_m]
  = \prod_{i=0}^{m-1} \left(1 - \frac{1}{2^w - i}\right)
  = \prod_{i=0}^{m-1} \frac{2^w - i - 1}{2^w - i}
  = \frac{2^w - m}{2^w}.
\]
Thus the probability that some fresh decoded candidate value equals the
hidden $Z$ is $m / 2^w \leq \qRO / 2^w \leq \qRO/2^k$, the last step by
$w \geq k$. This event is exactly the first $\badkey$ abort, because
$\KeyedTWOut_Z$ returns either $\bot$ or a unique candidate $Z$-value.
\end{proof}

\subsection{Leaf Tag Collision Bound}
\label{app:lem-leaf-collision}

\begin{lemma}[Leaf Tag Collision Bound]
\label{lem:leaf-collision}
Consider any \TW{} random oracle model game with an adversary
$\mathcal{B}$ that has direct access to $\ROflat$ and construction
computations that may generate final leaf transcripts---by encryption,
plaintext-computing, or verification computations, whether they accept or
reject. Let
$\badleaf$ be the event that two distinct construction-generated final
leaf transcript prefixes $X \neq X'$ in leaf tag output-point classes
receive the same $\tauleaf$-bit projection from $\RF(\cdot)$, and let
$\nleaf(\mathcal{B})$ be the execution-uniform count of distinct
construction-generated final leaf transcripts (\Cref{sec:resources}).
The first bound applies the collision case ($s = 2$, $\tau = \tauleaf$)
of \Cref{lem:adaptive-candidate} to a chronological leaf tag candidate
sequence; the second is a same-slot freshness argument in an all-reject
game.
\begin{enumerate}
\item \emph{General count.} Unconditionally---in particular when some
  or all keys are adversary-supplied, so direct queries may themselves
  land on leaf tag output points---
  \[
    \Pr[\badleaf] \;\leq\;
    \LeafColl_{\tauleaf}\bigl(\qRO(\mathcal{B})+\nleaf(\mathcal{B})\bigr)
    = \binom{\qRO(\mathcal{B})+\nleaf(\mathcal{B})}{2} \big/ 2^{\tauleaf}.
  \]
\item \emph{Hidden key, same slot.} Suppose the game has a hidden key,
  nonce-respecting encryption queries, and \emph{all-reject}
  verification: every non-replay verification submission performs its
  construction computation and returns reject to the adversary. (The
  game \Grej{} of \Cref{def:adjacent-hybrid-games} and the
  \textsf{INT-RUP} continuation in the proof of
  \Cref{lem:int-rup-root-guess} have this shape.) Say two final
  leaf transcripts \emph{share a slot} if they recover the same key,
  nonce, and leaf index $(K, U, j)$ (\Cref{cor:output-sep}). Let
  $\badleaf$ be the event that some construction-generated final leaf
  transcript receives the same $\tauleaf$-bit projection as a distinct
  final leaf transcript that shares its slot and is generated by an
  encryption computation, and let $\mathsf{hit}_{\mathsf{leaf}}$ be the
  event that a direct adversary query hits the bridged input of a
  construction-generated hidden-key final leaf transcript before that
  transcript is first created. Then, in such a game,
  \[
    \Pr[\badleaf \text{ before } \mathsf{hit}_{\mathsf{leaf}}]
    \;\leq\; \frac{\nleaf(\mathcal{B})}{2^{\tauleaf}}.
  \]
\end{enumerate}
\end{lemma}

\begin{proof}
Build a chronological leaf tag candidate sequence: a candidate is the
bridged input of a leaf tag output point ($\Ltag(j)$), appended the first
time it is reached---by a direct query in $\mathcal{B}$'s oracle phase or
by a construction computation generating a final leaf transcript---and
not appended again. By \Cref{lem:bridge-decode}, \Cref{cor:output-sep},
and \Cref{lem:transcript-inj}, two reaches of the same bridged input come
from the same final leaf transcript, so repeated evaluations add no
candidate---the lazy $\ROflat$ value is reused---and distinct final leaf
transcripts give distinct candidates. No candidate's $\ROflat$ value is
read before its append step, so the sequence meets the predictability and
unrevealed-value premises of \Cref{lem:adaptive-candidate}, whose
collision case with $s = 2$ and $\tau = \tauleaf$ bounds the probability
that two candidates share a $\tauleaf$-bit projection by
$\binom{L}{2} \cdot 2^{-\tauleaf}$ for a length-$L$ sequence. Every
$\badleaf$ pair is a pair of construction-generated candidates.

\emph{(1) General count.} Direct queries may decode to a leaf
tag class $\Ltag(j)$ and are appended as candidates; together with the
$\nleaf(\mathcal{B})$ construction-generated transcripts the sequence has
length at most $\qRO(\mathcal{B}) + \nleaf(\mathcal{B})$, so the case
with $L = \qRO(\mathcal{B}) + \nleaf(\mathcal{B})$ gives
$\Pr[\badleaf] \leq
\binom{\qRO(\mathcal{B})+\nleaf(\mathcal{B})}{2} \cdot 2^{-\tauleaf}
= \LeafColl_{\tauleaf}(\qRO(\mathcal{B})+\nleaf(\mathcal{B}))$. The
count charges every direct query and every construction-generated final
leaf transcript regardless of key provenance, so no assumption on the
keys is used.

\emph{(2) Hidden key, same slot.} Every $\badleaf$ pair is a pair of
construction-generated hidden-key transcripts. Nonce-respecting
encryption uses each nonce in at most one encryption query, and within
one encryption computation the leaf index separates final leaf
transcripts, so each slot contains at most one encryption-generated
final leaf transcript---which may be a transcript first created by an
earlier verification or plaintext-computing computation, when the
encryption's absorbed leaf body reproduces an earlier absorbed body.
Each construction-generated final leaf transcript therefore belongs to
at most one $\badleaf$ pair, and at most $\nleaf(\mathcal{B})$ pairs can
occur.

Each pair's collision is decided at the first creation of its
second-created member. At that moment the earlier member's
$\tauleaf$-bit projection is a function of the prior execution, while
the later member's projection is a fresh uniform lazy sample: by
Transcript Injectivity (\Cref{lem:transcript-inj}) two construction
reaches of the same bridged input come from the same final leaf
transcript, so no construction computation has read that bridged input
before the creation, and a direct query reading it before the creation
would decode under $\KeyedTWOut$ to the hidden key and its $\Ltag(j)$
class (\Cref{lem:bridge-decode})---precisely
$\mathsf{hit}_{\mathsf{leaf}}$. Hence before
$\mathsf{hit}_{\mathsf{leaf}}$, conditioned on the entire prior
execution, each pair collides on its $\tauleaf$-bit projections with
probability $2^{-\tauleaf}$, and a union bound over the at most
$\nleaf(\mathcal{B})$ pairs gives the claim. No assumption on what the
visible interface exposes is needed: a projection revealed after its
sampling cannot alter a collision already decided. By Leaf-Transcript
Recovery (\Cref{lem:leaf-recovery}), on $\neg\badleaf$ equal flattened
final-root inputs of an encryption computation and any construction
computation identify the same leaf transcript sequence.
\end{proof}

\subsection{Encryption Transcript Marginal Uniformity}
\label{app:enc-transcript-marginal-uniformity}

\begin{lemma}[Encryption Transcript Marginal Uniformity]
\label{lem:enc-transcript-marginal-uniformity}
For any fixed nonce-respecting encryption query in the \TW{}
random oracle experiment, let $\mathsf{hit}_{\mathsf{fresh}}$ be the
event that, before this query reaches some counted transcript point,
a prior direct adversary query to $\ROflat$ has already hit the
bridged input of that hidden-key point. On
$\neg\mathsf{hit}_{\mathsf{fresh}}$, every transcript output of the
query that contributes to a returned ciphertext block, to an internal
leaf tag, or to the final root tag is either fresh when this query first
generates its bridged input, or reuses the lazy-table value sampled
uniformly when that input was first generated by an earlier construction
computation. The counted output points first generated by the fixed
query are pairwise distinct.
\end{lemma}

\begin{proof}
Work on an execution outside $\mathsf{hit}_{\mathsf{fresh}}$ for the
fixed query. This rules out prior adversary reads of the hidden-key
construction inputs considered below. A previous construction computation
may nevertheless have generated one of these inputs, for example through
a non-replay verification query using the same nonce. Such a lookup
samples the lazy-table value uniformly when it first occurs; if the fixed
encryption query later reaches the same bridged input, it reuses that
same value. Thus it remains only to show that the counted output points
newly generated by the fixed encryption query are distinct from one
another and from previous nonce-respecting encryption queries. Conditional
information leaked by earlier verification comparisons of a final root
tag is accounted for separately by \Cref{lem:ver-first-root-guess}, not
by this lemma.

\paragraph{First root keystream block.}
The first root keystream block, when present, is produced by the root
duplex after the initialization call $\prfxroot \concat K
\concat U$ and any AD absorption (elided when $A = \varepsilon$, in
which case the \siginit call itself emits the first root keystream
block). Since the nonce-respecting adversary never repeats
a nonce, no previous encryption query can have used the same root
initialization transcript. Within the current query,
\Cref{lem:transcript-inj} separates initialization, AD, message, and
aggregation calls by their suffixes and by the duplex-call encoding,
so the first root keystream transcript is distinct from the other
counted output points generated by this query. If the requested
keystream length is zero, no ciphertext block is returned at that point
and there is nothing to prove for that output.

\paragraph{First leaf keystream block.}
For each leaf chunk $j \geq 1$, the first leaf keystream block is
produced after the initialization call $\prfxleaf \concat
K \concat U \concat \langle j \rangle_8$. The prefix separates leaves
from the root, the nonce separates this encryption query from all
previous encryption queries, and $\langle j \rangle_8$ separates
leaves within this query. Hence every first leaf keystream transcript
newly generated by the query is distinct.

\paragraph{Later message outputs.}
Argue inductively within each root or leaf message stream. Suppose
the current message-output transcript is either first generated by this
query or reuses a value generated by an earlier construction computation.
If its output is used as a keystream block, write its $\ROflat$ value as
$Z_i$. Since the adaptive plaintext block $M_i$ is a deterministic
function of the prior view,
\[
  C_i = M_i \oplus Z_i
\]
is now fixed. \TW{} then absorbs $C_i$ with the appropriate message
phase suffix ($\sigmsgmore$ or
$\sigmsglast$) before requesting the next keystream block. Once $C_i$ is fixed, the next transcript is
fixed; \Cref{lem:transcript-inj}'s injectivity and domain separation
imply it has not appeared earlier unless the same duplex-call prefix
has appeared earlier. Nonce uniqueness excludes this across encryption
queries. By \Cref{cor:output-sep}, the next message-keystream output
point can coincide with a within-query earlier output point only if
both belong to the same output-point class on the same duplex
instance (root or leaf-$j$ with matching $j$) and have the same
duplex-call sequence; the current message-stream prefix is strictly
longer than every earlier prefix in that stream, so this is excluded.
Across distinct instances within the same query (root vs leaf, or
two leaves with distinct $j$), the prefixes already differ in the
first call's $\sigma$ by $\prfxroot /
\prfxleaf$ or $\langle j \rangle_8$ respectively. Each
later message-output transcript that emits keystream is therefore
distinct from the other points generated by this query, unless it is a
reuse of a point generated by an earlier construction computation. The
terminal message-output transcript of a leaf emits
the leaf tag rather than a following keystream block; the same
separation argument applies to that transcript, so each internal leaf
tag newly generated by the query is sampled at a distinct point. Empty
inputs still induce a duplex call, but with zero keystream output they
contribute no ciphertext block.

\paragraph{Final root tag transcript.}
The final root tag is sampled at the root tag output point selected by
$\|C\|$: the $\Rtag$ point after the root duplex absorbs the
aggregation transcript built from the leaf tags and $\langle n
\rangle_8$ when $n \geq 1$, and the $\Rmsg^+$ point closing the root
message phase when $n = 0$ (\Cref{sec:enc}). When $n \geq 1$, the leaf
tags are themselves $\ROflat$ outputs on final leaf transcripts;
condition on their sampled values, on the sampled
ciphertext, and on the adversary-visible prior view and $\ROflat$
table state, so that the root aggregation transcript is fixed. When
$n = 0$ there is no aggregation transcript; condition on the sampled
ciphertext and the adversary-visible prior view alone. By
\Cref{cor:output-sep}, the selected root tag output-point class is
separated from every other output-point class in the same query ---
including root initialization (suffix \siginit), AD-phase outputs,
the remaining message phase outputs, leaf-side output points (separated
by $\prfxleaf$), and, in the $n \geq 1$ case, earlier \Ragg-prefix
output points (separated by the \sigaggmore vs
\sigagglast suffix). Even if two encryption queries realize
the same aggregation transcript (or, for $n = 0$, the same root
ciphertext), the full root transcript also
contains the root initialization $\prfxroot \concat K
\concat U$, and nonce uniqueness separates that initialization from
all previous encryption queries. The final root tag transcript is
therefore either newly generated at a distinct point of the fixed query,
or is a reuse of a value first sampled by an earlier construction
computation. In the newly generated case, its $\tauroot$-bit $\ROflat$
output is a fresh lazy-table sample; in the reuse case, it is the same
uniform sample drawn when the bridged input was first generated.
\end{proof}

\subsection{Mu-Ind Hybrid Games and Verification Accounting}
\label{app:mu-verification-coupling}

The direct mu-ind adjacent-hybrid proof bounds one user-replacement step
of the hybrid through an explicit game sequence, defined next, using the
generic key-test and leaf-collision accounting lemmas together with the
verification-specific facts below. We use the root tag output point and
final root transcript classes of \Cref{def:root-tag-output-point}; in
particular, computations on a common ciphertext share the same root tag
output point because $\leaves(C)$ is fixed by $\|C\|$.

\begin{definition}[Adjacent-Hybrid Games]
\label{def:adjacent-hybrid-games}
Fix a flat-RO mu-ind adversary $\mathcal{B}$, the pre-sampled keys
$K_1, \dots, K_D$ of \Cref{sec:mu-setting}, and a user index
$d \in \{1, \dots, D\}$. On the branch where two sampled keys collide,
every game below outputs the same fixed bit; the remaining clauses
describe the branch with pairwise distinct keys. Every game maintains
one lazy table $\Theta$ implementing $\ROflat$: $\Theta[Y]$ samples a
fresh uniform infinite string at the first access to $Y$ and returns
the stored string afterwards---equivalently, $\Theta$ is sampled in
full at the outset and read on demand. Direct adversary queries
$(Y, \ell)$ are answered by $\lfloor \Theta[Y] \rfloor_\ell$; the
\emph{environment} answers users $e < d$ by $(\Rand, \Replay)$ and
users $e > d$ by $(\Enc_{K_e}^\RO, \Ver_{K_e}^\RO)$, reading
construction outputs from $\Theta$; and the $\mathsf{New}$ interface
and the validity and nonce-respecting checks of \Cref{sec:mu-setting}
are common to every game. Every game maintains the flags $\badkey$,
$\badleaf$, and $\forge$, initialized false; no oracle answer or branch
depends on a flag, and $\badkey$ is set whenever a direct query input
decodes under $\KeyedTWOut$ to $K_d$ and a counted output-point class
of \Cref{tab:keyed-twout}.
\begin{itemize}
\item Games \Greal{} and \Grej{} answer user $d$ by the shared oracle
  code of \Cref{alg:adjacent-hybrid-games} and differ exactly in the
  marked return statement: a non-replay verification whose comparison
  succeeds sets $\forge$ and returns accept in \Greal{} and reject in
  \Grej{}.
\item Game \Gideal{} answers user-$d$ encryption queries with fresh
  uniform byte strings of length $\|M\| + \tauroot/8$, recorded in
  $W$, and user-$d$ verification by the replay check alone; it performs
  no user-$d$ construction computation and sets neither $\badleaf$ nor
  $\forge$. \Gideal{} is the hybrid $H_d^\star$ of
  \Cref{lem:mu-adjacent-hybrid}, instrumented with $\badkey$.
\end{itemize}
The flag $\badleaf$ is set, in \Greal{} and \Grej{}, when a user-$d$
construction computation first creates a final leaf transcript whose
$\tauleaf$-bit projection equals that of a distinct, already created
final leaf transcript sharing its $(K_d, U, j)$ slot, at least one of
the two generated by an encryption computation---the same-slot event of
\Cref{lem:leaf-collision}, part~2, for user $d$. It guards no code
difference and is used only in the accounting.
\end{definition}

\begin{algorithm}[t]
  \caption{Shared user-$d$ oracle code of the games \Greal{} and
    \Grej{} (\Cref{def:adjacent-hybrid-games}). $\Theta$ is the shared
    lazy $\ROflat$ table, read at bridged transcript prefixes, and $W$
    the user-$d$ replay record. Inputs failing the validity or
    nonce-respecting checks of \Cref{sec:mu-setting} are rejected
    beforehand, as in every game. The two games differ exactly in the
    marked lines~\ref{line:ret-real} and~\ref{line:ret-rej}.}
  \label{alg:adjacent-hybrid-games}
  \begin{algorithmic}[1]
    \Procedure{RO}{$Y$, $\ell$} \Comment{direct adversary query}
      \State set $\badkey$ if $\KeyedTWOut(Y) = (K_d, \mathcal{C})$ for
        a counted class $\mathcal{C}$ (\Cref{tab:keyed-twout})
      \State \Return $\lfloor \Theta[Y] \rfloor_{\ell}$
    \EndProcedure
    \Statex
    \Procedure{Enc}{$U$, $A$, $M$}
      \State run $\Enc_{K_d}^{\RO}(U, A, M)$, reading each requested
        duplex output as the corresponding projection of $\Theta$ at
        its bridged transcript prefix (\Cref{sec:ro-model}); let the
        result be $C \concat T$
      \State set $\badleaf$ if a final leaf transcript first created
        here completes a same-slot collision
        (\Cref{def:adjacent-hybrid-games})
      \State record $(U, A, C \concat T)$ in $W$; \Return $C \concat T$
    \EndProcedure
    \Statex
    \Procedure{Vf}{$U$, $A$, $C \concat T$}
      \If{$(U, A, C \concat T) \in W$} \Return accept \EndIf
      \State determine the verification transcript of $(K_d, U, A, C)$;
        read from $\Theta$ the $\tauleaf$-bit leaf tags at its final
        leaf transcripts, setting $\badleaf$ as in \textsc{Enc}, and
        nothing at its stream output points
      \State $T' \gets \lfloor \Theta[\flatmap(R)] \rfloor_{\tauroot}$
        at the root tag output point $R$ of the computation
        (\Cref{def:root-tag-output-point})
      \If{$T' = T$}
        \State $\forge \gets \mathsf{true}$
        \State \Return accept \Comment{in \Greal{} only}
          \label{line:ret-real}
        \State \Return reject \Comment{in \Grej{} only}
          \label{line:ret-rej}
      \EndIf
      \State \Return reject
    \EndProcedure
  \end{algorithmic}
\end{algorithm}

\begin{lemma}[Environment Neutrality]
\label{lem:env-neutrality}
In each game of \Cref{def:adjacent-hybrid-games}, on the branch with
pairwise distinct keys, the environment's construction-internal
$\ROflat$ lookups---made while simulating $\Enc_{K_e}^\RO$ and
$\Ver_{K_e}^\RO$ for users $e > d$---neither set a user-$d$ flag nor
generate the bridged input of a user-$d$ counted transcript point.
\end{lemma}

\begin{proof}
Whenever such a lookup lies in a counted keyed transcript class,
\Cref{lem:bridge-decode} decodes its candidate key as $K_e$, which on
the distinct-keys branch differs from $K_d$; a lookup outside the
counted keyed language is irrelevant to the flags. By the key-field
disjointness of \Cref{lem:key-renaming}, the environment's addresses
are disjoint from user $d$'s. The flags are set only by direct
adversary queries decoding to $K_d$ ($\badkey$), by user-$d$ final leaf
transcript creations ($\badleaf$), and by user-$d$ verification
comparisons ($\forge$), none of which the environment performs. Ideal
users $e < d$ perform no construction lookups at all.
\end{proof}

\begin{lemma}[Verification Visible-Prefix Key Uniformity]
\label{lem:ver-key-uniformity}
In the stopped single-key \TW{} random oracle encryption/verification
interaction, or its \textsf{INT-RUP} extension with a
plaintext-computing oracle, with hidden key $K$, direct access to $\ROflat$, and
immediate abort on $\badkey$, order the direct $\ROflat$ queries. For a
direct query position $i$, let $S_i$
be the set of distinct candidate keys decoded by $\KeyedTWOut$ from
earlier direct-query inputs. Condition on any positive-probability
visible execution prefix immediately before query $i$, on no earlier
$\badkey$, and on $S_i$. Then the hidden key $K$
is uniform on $\bits^k \setminus S_i$. Consequently the uniformity
invariant required by \Cref{lem:key-tests} holds in the stopped
interaction.
\end{lemma}

\begin{proof}
The visible prefix contains the adversary's oracle queries and replies:
direct $\ROflat$ outputs, encryption outputs and the resulting replay
table, plaintext-computing replies when that oracle is present,
verification bits, and the adversary state determined by these
values. Fix such a prefix and two keys $K_0, K_1 \notin S_i$.

Let $\Phi_{0 \to 1} = \Phi_{K_0 \to K_1}$ be the key-renaming map of
\Cref{lem:key-renaming} and $\Phi_{1 \to 0}$ its inverse. By that lemma,
$\flatmap \circ \Phi_{0 \to 1} \circ \flatmap^{-1}$ is a bijection
$\mathcal{Y}_{K_0} \to \mathcal{Y}_{K_1}$ between the bridged hidden-key
construction addresses for $K_0$ and for $K_1$, preserving output-point
classes and equality relations among addresses.

This bijection is applied online, not to a pre-existing fixed address
set. Whenever a construction computation under $K_0$ would first query
a hidden-key address $Y$, the coupled computation under $K_1$ queries
$\Phi_{0 \to 1}(Y)$ and receives the same lazy-table value; repeated
queries are coupled consistently through the same map. If the queried
address is a root aggregation address containing leaf tag bytes, those
bytes are already coupled equal because the corresponding leaf tag
lookups were either repeated or first sampled earlier under the same
renaming rule. Thus adaptively generated message transcripts and root
aggregation transcripts are renamed step by step with the same visible
ciphertext, leaf tag, and root tag bytes.

Prior direct-query inputs are not renamed: they are adversary-chosen
visible strings. Since $K_0, K_1 \notin S_i$, neither
$\mathcal{Y}_{K_0}$ nor $\mathcal{Y}_{K_1}$ contains a prior direct-query
input (\Cref{lem:key-renaming}, direct-query avoidance). Direct-query
lazy-table values are therefore coupled identically and remain disjoint
from the renamed hidden
construction addresses.

Use the same adversary coins, the same direct-query lazy-table values,
and the online renamed construction-internal lazy-table values.
Encryption outputs have the same distribution under the coupling,
because every construction lookup that determines ciphertext, leaf
tags, or the final root tag is paired with a distinct renamed lookup
carrying the same sampled value. Plaintext-computing computations are
renamed in the same online way: a plaintext reply is the submitted
ciphertext block XORed with a coupled keystream projection at a stream
output point (\Cref{tab:keyed-twout}), so it carries the same visible
bytes under both keys, and the recomputed tag is not returned. Replay
verification decisions depend
only on the visible replay table. Non-replay verification computations
are renamed in the same online way as encryption computations, and the
final root tag comparison therefore gives the same accept/reject bit.
Verification does not reveal plaintext on rejection or acceptance; it
reveals only that bit.

Thus every visible prefix with hidden key $K_0 \notin S_i$ has the
same conditional probability as the corresponding execution with hidden
key $K_1 \notin S_i$. Since $K$ is initially uniform and the condition
``no earlier $\badkey$'' removes exactly the keys in
$S_i$, all keys in $\bits^k \setminus S_i$ remain equally likely.
\end{proof}

\paragraph{Foreign-key key-test claim.}
In each game of \Cref{def:adjacent-hybrid-games}, work in the original
experiment, where the key-collision branch outputs a fixed bit rather
than conditioning on pairwise distinct keys. Then
\[
  \Pr[\badkey] \leq \KeyGuess_k(\qRO(\mathcal{B})).
\]
Fix the foreign keys and leave their transcript addresses unchanged.
Before the first user-$d$ key hit, the visible prefix is invariant under
renaming $K_d$ among keys that have not already been tested and are not
equal to a fixed foreign key, by the same online coupling as
\Cref{lem:ver-key-uniformity} and the key-field disjointness of
\Cref{lem:key-renaming}. For each fresh candidate key tested by one of
$\mathcal{B}$'s direct $\ROflat$ queries, the contributing event is
contained in $\{K_d \text{ equals that candidate}\}$ and costs at most
$2^{-k}$ in the original pre-sampled-key experiment; a candidate equal
to a fixed foreign key contributes only on the fixed-output
key-collision branch and therefore contributes nothing to
$\badkey$. A union bound over at most $\qRO(\mathcal{B})$ direct
queries gives the claim.

\begin{lemma}[Verification-First Root Guessing]
\label{lem:ver-first-root-guess}
In the game \Grej{} of \Cref{def:adjacent-hybrid-games}, stopped at the
first occurrence of $\badkey \cup \badleaf$---with $\badleaf$ the
same-slot event of \Cref{lem:leaf-collision}, part~2, for user
$d$---let $\mathsf{win}_{\mathsf{vf}}$ be the event that some
non-replay user-$d$ verification comparison succeeds within a
verification-first final root transcript class. Then
\[
  \Pr[\mathsf{win}_{\mathsf{vf}}]
  \leq \qv^{(d)} / 2^{\tauroot},
\]
where $\qv^{(d)}$ is the user-$d$ non-replay verification cap of
\Cref{sec:mu-setting}.
\end{lemma}

\begin{proof}
Throughout, a \emph{non-replay verification submission} means a valid
non-replay user-$d$ verification query; by the validity convention, an
invalid verification query returns $0$ and performs no construction
computation, so it realizes no bridged final root input, belongs to no
final root transcript class, and contributes to no $G_{\mathcal{C}}$
below.

In \Grej{}, every non-replay submission performs its construction
computation and returns reject (\Cref{alg:adjacent-hybrid-games}), so
no comparison outcome is visible to the adversary, and each class's
root tag value is read from the lazy table only by the class's own
comparisons and by any later encryption computation reaching the class.

Fix a verification-first final root transcript class $\mathcal{C}$ and
follow the stopped execution until the first encryption query
generating the same class, if any (the cutoff). Let $G_{\mathcal{C}}$
be the set of distinct candidate tags submitted by non-replay
verification queries in class $\mathcal{C}$ before that cutoff;
repeated tags only shrink this set. The class's correct root tag
$U_{\mathcal{C}}$ is the $\tauroot$-bit projection of $\ROflat$ at the
class's bridged root tag input. Before the stop, no direct query has
read this hidden-key input: such a query would be decoded by
$\KeyedTWOut$ to the hidden key and would set $\badkey$; the
environment reads no user-$d$ address (\Cref{lem:env-neutrality}).
Before the cutoff, no encryption query has revealed $U_{\mathcal{C}}$
as an encryption tag. (When the class's root tag point is $\Rmsg^+$,
i.e.\ $n = 0$, an encryption query on a longer message that shares
chunk $0$ may touch this same address while closing its chunk-$0$
message phase, but it requests zero output there and continues into its
aggregation phase, so it does not reveal $U_{\mathcal{C}}$; its own tag
is emitted at a later $\Rtag$ point.) The comparisons against
$U_{\mathcal{C}}$ return the constant reject, so the visible execution
and the set $G_{\mathcal{C}}$ are functions only of the adversary's
coins, visible oracle replies, and lazy-table values at inputs other
than this class's root tag address. By the lazy-sampling order
exchange, conditioned on those values, $U_{\mathcal{C}}$ may
equivalently be sampled after $G_{\mathcal{C}}$ is fixed, and is
uniform independently of $G_{\mathcal{C}}$. The probability that some
pre-cutoff comparison in class $\mathcal{C}$ succeeds---that
$U_{\mathcal{C}} \in G_{\mathcal{C}}$---is therefore at most
$|G_{\mathcal{C}}|/2^{\tauroot}$.

After the cutoff, if it exists, class $\mathcal{C}$ admits no
successful non-replay comparison. Let the cutoff encryption query
return $(U_e,A_e,C_e,T_e)$. Any later verification submission in the
same final root transcript class has, by the final-root determinacy of
\Cref{lem:transcript-inj} and Leaf-Transcript Recovery
(\Cref{lem:leaf-recovery}) under the absence of $\badleaf$---which
supplies the recovery hypothesis because the cutoff computation is an
encryption computation---the same $(U,A,C)$ as that encryption
computation. If the submitted tag is $T_e$, the tuple is a recorded
encryption tuple and the query is answered by the replay check before
any comparison; if the submitted tag is not $T_e$, the comparison
fails.

Each non-replay verification submission deterministically reconstructs
one final root transcript at the root tag output point, so it
contributes its submitted tag to at most one set $G_{\mathcal{C}}$, and
only when its class is verification-first and the submission precedes
that class's cutoff. Thus
$\sum_{\mathcal{C}} |G_{\mathcal{C}}| \leq \qv^{(d)}$, and a union
bound over the verification-first classes gives the claim.
\end{proof}

\begin{lemma}[\textsf{INT-RUP} Root Guessing]
\label{lem:int-rup-root-guess}
In the stopped single-key \TW{} random oracle
encryption/plaintext-computing/verification interaction before
$\badkey \cup \badleaf$---with $\badleaf$ the same-slot event of
\Cref{lem:leaf-collision}, part~2---let $\forge$ be the event that some valid
non-replay verification query accepts. Then
\[
  \Pr[\forge \text{ before } \badkey \cup \badleaf]
  \leq \qv(\mathcal{B}) / 2^{\tauroot}.
\]
\end{lemma}

\begin{proof}
Throughout, a \emph{non-replay verification submission} means a valid
non-replay verification query; invalid verification queries return $0$,
perform no construction computation, realize no bridged final root
input, and contribute to no set of guessed tags.

Work in the game stopped at the first occurrence of
$\badkey$ or $\badleaf$. A final root transcript class is
\emph{tag-revealed} once an encryption query has returned the root tag
for that class. If a non-replay verification submission belongs to a
tag-revealed class, compare it with an encryption computation that
revealed the class. The two computations' final-root transcript
prefixes are well-formed typed transcript prefixes with the same
bridged final root input; by
\Cref{lem:bridge-decode}, they have the same native flattened input.
Then final-root determinacy (\Cref{lem:transcript-inj}) and
Leaf-Transcript Recovery (\Cref{lem:leaf-recovery}) under
$\neg\badleaf$---which supplies the recovery hypothesis because the
comparand is an encryption computation---recover the same $(U,A,C)$. If the submitted tag is the
revealed tag, the tuple is a replay-table hit; if it is any other tag,
the final comparison rejects. Thus any accepting non-replay
verification submission must lie in a class whose root tag has not yet
been revealed by encryption.

For tag-unrevealed classes, define an all-reject continuation:
encryption and plaintext-computing queries are answered normally,
replay verification queries accept normally, and each non-replay
verification submission performs the same construction computation but
returns reject to the adversary. Plaintext-computing queries may
generate the same stream and leaf tag transcripts as decryption, but
they do not reveal the final root tag value. Equivalently, the final
root tag sample for a tag-unrevealed class can be deferred until it is
first needed for an encryption answer or a verification comparison;
this does not change any plaintext reply. By Output-Point Separation
(\Cref{cor:output-sep}) and the keystream schedule recorded at
\Cref{tab:keyed-twout}, the keystream values used to compute returned
plaintext blocks are projections at stream output points, not at the
final root tag output point; in particular the chunk-$0$ $\Rmsg^+$ call
of an $n \geq 1$ computation requests zero output, so no plaintext
reply exposes an $\Rmsg^+$ tag value.

Fix a final root transcript class $\mathcal{C}$ whose first
construction generation is not by an encryption query; for a class
first generated by an encryption query the tag is revealed from that
generation onward, so every verification submission in it is covered by
the previous paragraph. Follow the all-reject path until the first
encryption query that reveals the class, if any (the cutoff; any such
query is later than the class's first generation). Let
$G_{\mathcal{C}}$ be the set of distinct
tags submitted by non-replay verification queries in class
$\mathcal{C}$ before that cutoff. Before the cutoff, no encryption query
has returned the class tag and no direct query has read the hidden-key
root tag point without triggering $\badkey$. Conditioned on the
all-reject path and on all lazy-table values except this final root tag
projection, the correct tag for $\mathcal{C}$ is uniform in
$\bits^{\tauroot}$ and independent of $G_{\mathcal{C}}$. The
probability that any pre-cutoff submitted tag in $G_{\mathcal{C}}$
equals the deferred class tag is therefore at most
$|G_{\mathcal{C}}|/2^{\tauroot}$.

Consider the event that the first accepting non-replay verification
submission in the stopped game belongs to class $\mathcal{C}$. Up to
that global first accept, the adversary's visible answers match the
all-reject continuation: every earlier non-replay submission rejects,
replay and plaintext-computing queries are answered identically in both
executions, and encryption answers are unchanged. The construction
computations also agree at every address except that the continuation
defers the final root tag value for class $\mathcal{C}$ until the
comparison is needed. Therefore the first accepting non-replay
submission can belong to class $\mathcal{C}$ only before the cutoff,
and only if its submitted tag equals the deferred class tag, which has
probability at most $|G_{\mathcal{C}}|/2^{\tauroot}$.

The execution-uniform count $\qv(\mathcal{B})$ applies to the
all-reject continuation. Each valid non-replay verification query
contributes its submitted tag
to at most one such set $G_{\mathcal{C}}$, and replay-table hits do not
contribute. Thus
$\sum_{\mathcal{C}} |G_{\mathcal{C}}| \leq \qv(\mathcal{B})$, and a
union bound over the classes not first generated by an encryption
query, each with its pre-cutoff guess set $G_{\mathcal{C}}$, gives the
claim.
\end{proof}

\subsection{Mu-Ind Adjacent-Hybrid Lemma}

The multi-user indistinguishability theorem (\Cref{thm:mu-ind}) telescopes a
hybrid over the $D$ users, replacing one real user by its ideal
$(\Rand,\Replay)$ interface at each step. The following lemma bounds one such
step through the game sequence of \Cref{def:adjacent-hybrid-games}: the
instrumented real game \Greal{}, the all-reject game \Grej{} differing from
it in a single return statement, and the ideal game \Gideal{}, with the
divergence events set as the flags $\badkey$, $\badleaf$, and $\forge$ by
game code.

\begin{lemma}[Mu-Ind Adjacent-Hybrid Bound]
\label{lem:mu-adjacent-hybrid}
Let $\mathcal{B}$ be a flat-RO mu-ind adversary in the pre-sampled-key
formulation of \Cref{sec:mu-setting}, with keys
$K_1,\dots,K_D$ sampled at the start of the experiment. Fix
$d \in \{1,\dots,D\}$. Let $H_{d-1}^\star$ and $H_d^\star$ be the
adjacent stopped hybrids that output a fixed bit if any two sampled keys
collide, and otherwise run $\mathcal{B}$ with one shared lazy
$\ROflat$ table as follows: users $e < d$ are answered by
$(\Rand,\Replay)$, users $e > d$ are answered by
$(\Enc_{K_e}^\RO,\Ver_{K_e}^\RO)$, and user $d$ is answered by
$(\Enc_{K_d}^\RO,\Ver_{K_d}^\RO)$ in $H_{d-1}^\star$ and by
$(\Rand,\Replay)$ in $H_d^\star$. If $\nleaf^{(d)}$ and $\qv^{(d)}$
are the user-$d$ resource caps, then
\[
  \bigl|\Pr[\mathcal{B}^{H_{d-1}^\star} \Rightarrow 1]
       - \Pr[\mathcal{B}^{H_d^\star} \Rightarrow 1]\bigr|
  \leq
  \KeyGuess_k(\qRO(\mathcal{B}))
  + \frac{\nleaf^{(d)}}{2^{\tauleaf}}
  + \frac{\qv^{(d)}}{2^{\tauroot}}.
\]
\end{lemma}

\begin{proof}
The fixed-output branch on a key collision is identical in
$H_{d-1}^\star$, $H_d^\star$, and the games of
\Cref{def:adjacent-hybrid-games}, so only the branch with pairwise
distinct keys contributes; work on that branch, with the keys fixed,
throughout. The proof identifies $H_{d-1}^\star$ with \Greal{}
(Claim~1), relates \Greal{} to \Grej{} pointwise (Claim~2) and \Grej{}
to $H_d^\star = \Gideal{}$ in distribution (Claim~3), and assembles the
bound by a single first-occurrence decomposition in \Grej{} (Claim~4).

\paragraph{Claim 1: \Greal{} realizes $H_{d-1}^\star$.}
The user-$d$ oracles of \Cref{alg:adjacent-hybrid-games} implement
$(\Enc_{K_d}^\RO, \Ver_{K_d}^\RO)$ with three distribution-preserving
refactors. First, reading every construction output as a projection of
the shared table $\Theta$ at its bridged transcript prefix is the lazy
implementation of $\ROflat$ (\Cref{sec:ro-model}); whether $\Theta$ is
sampled in full at the outset or entry by entry on first access is
immaterial. Second, \textsc{Vf} reads $\Theta$ only at leaf tag and
root tag output points. This is sound because decryption absorbs the
submitted ciphertext, so the verification transcript is determined by
$(K_d, U, A, C)$ independently of any stream output
(\Cref{lem:duplex-flat}, \Cref{sec:dec}), and verification discards the
recovered plaintext, so the keystream projections a real
$\Ver_{K_d}^\RO$ computation reads influence neither a visible answer
nor the value of any other table entry; leaving those entries unread
until some later computation first accesses them leaves the joint
distribution of all oracle answers unchanged. Third, the flags are
instrumentation: no oracle answer or branch depends on them. Hence
$\Pr[\mathcal{B}^{\Greal} \Rightarrow 1]
= \Pr[\mathcal{B}^{H_{d-1}^\star} \Rightarrow 1]$.

\paragraph{Claim 2: \Greal{} and \Grej{} agree until $\forge$.}
Run the two games on the same adversary coins, environment coins, and
table realization $\Theta$. The games share all code, including the
comparison that sets $\forge$, and differ in exactly one statement, the
marked return line of \Cref{alg:adjacent-hybrid-games}, which executes
only when $\forge$ has just been set. The two executions are therefore
identical---every oracle answer, table access, record in $W$, and flag
update---up to and including the query at which $\forge$ is first set,
and the first differing value is that query's returned bit. In
particular, on executions in which \Grej{} never sets $\forge$, the two
runs coincide, outputs included, and the flags $\badkey$ and $\badleaf$
are set at the same queries in both games up to the first $\forge$.

\paragraph{Claim 3: \Grej{} agrees with $H_d^\star$ until $\badkey$.}
Couple $H_d^\star = \Gideal{}$ to \Grej{} as follows: run \Grej{} on
the same adversary and environment coins, give the \Gideal{} coordinate
the same visible oracle answers for every query answered while
$\badkey$ is unset, and let the \Gideal{} coordinate continue
independently according to its own law from the query that sets
$\badkey$ onward. We claim the resulting \Gideal{} coordinate is distributed
exactly as \Gideal{}. The answer rules coincide query by query except
for user-$d$ encryption: non-replay user-$d$ verification returns
reject in \Grej{} and under $\Replay$ alike, and replay queries are
answered from the record $W$, which the coupling keeps equal in both
coordinates; the environment runs identical code in both games; and
direct queries read table entries that agree in both games before
$\badkey$---an entry written by a user-$d$ construction computation is
read by a direct query only if that query decodes under $\KeyedTWOut$
to $K_d$ and a counted class (\Cref{lem:key-renaming}), which sets
$\badkey$, while entries written by the environment are produced
identically in both games (\Cref{lem:env-neutrality}).

It remains to show that, before $\badkey$, each user-$d$ encryption
answer of \Grej{} is, conditioned on the entire prior visible
transcript, a fresh uniform byte string of length $\|M\| +
\tauroot/8$---the law of $\Rand$. Fix such a query $(U, A, M)$. Before
$\badkey$, no direct query has read the bridged input of a user-$d$
hidden-key counted point, so the hypothesis
$\neg\mathsf{hit}_{\mathsf{fresh}}$ of Encryption Transcript Marginal
Uniformity (\Cref{lem:enc-transcript-marginal-uniformity}) holds, and
every counted output point first generated by this query is pairwise
distinct and receives a fresh uniform sample. The stream output points
are always first generated by this query: user-$d$ verification
computations read no stream points (\Cref{alg:adjacent-hybrid-games}),
nonce-respecting encryption and Transcript Injectivity
(\Cref{lem:transcript-inj}) separate them from every earlier user-$d$
encryption computation, and the environment writes only foreign-key
addresses (\Cref{lem:env-neutrality}). The ciphertext body $C$ is
therefore a fresh uniform string. The internal leaf tags are either
fresh samples or reuse entries written by user-$d$ verification
computations; they are never returned and influence the answer only
through the root tag address. The root tag $T$ is read at the
computation's root tag output point (\Cref{def:root-tag-output-point}),
an address distinct from every stream point of the same query by
Output-Point Separation (\Cref{cor:output-sep}). If that entry is first
generated here, $T$ is a fresh uniform sample. Otherwise the entry was
written by an earlier user-$d$ verification computation---before
$\badkey$, the only other writer of user-$d$ addresses---whose reads of
it produced only the constant reject answers of \Grej{}; since no
oracle answer or branch depends on the flags, the visible transcript is
a function of the coins and of table entries other than this one, and
conditioned on the visible transcript the entry remains uniform. In
either case $(C, T)$ has the $\Rand$ law, and at most one user-$d$
encryption reaches any root tag address, since two would share the
nonce $U$, contradicting nonce-respecting encryption. By induction over
the query sequence, the coupled \Gideal{} coordinate has exactly the
law of $H_d^\star$, and its visible transcript agrees with \Grej{}'s
while $\badkey$ is unset.

\paragraph{Claim 4: assembly.}
Place the three games on one probability space: shared adversary and
environment coins and one table realization $\Theta$ drive \Greal{} and
\Grej{} pointwise, and the $H_d^\star$ coordinate is coupled to \Grej{}
as in Claim~3. Until the first occurrence of $\forge \vee \badkey$ in
\Grej{}, the three visible transcripts coincide: \Greal{}'s by Claim~2,
since $\forge$ has not yet been set, and $H_d^\star$'s by Claim~3,
since $\badkey$ has not. When neither flag is ever set, the final
outputs of $\mathcal{B}$ coincide as well, so the outputs can differ
only on the event $\forge \vee \badkey$. With Claim~1 identifying
$\Pr[\mathcal{B}^{H_{d-1}^\star} \Rightarrow 1] =
\Pr[\mathcal{B}^{\Greal} \Rightarrow 1]$ and Claim~3 identifying
$\Pr[\mathcal{B}^{H_d^\star} \Rightarrow 1]$ with the coupled
coordinate's output probability,
\[
  \bigl|\Pr[\mathcal{B}^{H_{d-1}^\star} \Rightarrow 1]
       - \Pr[\mathcal{B}^{H_d^\star} \Rightarrow 1]\bigr|
  \;\leq\;
  \Pr_{\Grej}[\forge \vee \badkey].
\]
Decompose by first occurrence, inserting $\badleaf$---which guards no
code difference and enters only this accounting:
\[
\begin{aligned}
  \Pr_{\Grej}[\forge \vee \badkey]
  \;\leq\;
  \Pr[\badkey]
  &+ \Pr[\badleaf \text{ before } \badkey \vee \forge] \\
  &+ \Pr[\forge \text{ before } \badkey \vee \badleaf].
\end{aligned}
\]
The three terms are bounded in \Grej{}. For the first, the foreign-key
key-test claim above gives
$\Pr[\badkey] \leq \KeyGuess_k(\qRO(\mathcal{B}))$. For the second,
\Grej{} restricted to user $d$ is a hidden-key game with
nonce-respecting encryption and all-reject verification; before
$\badkey$ no direct query has read a user-$d$ hidden final leaf
point---such a query would decode to $K_d$ and its $\Ltag(j)$ class and
set $\badkey$---and the environment creates no user-$d$ leaf transcript
(\Cref{lem:env-neutrality}), so the same-slot case of
\Cref{lem:leaf-collision} gives
$\Pr[\badleaf \text{ before } \badkey \vee \forge]
\leq \nleaf^{(d)}/2^{\tauleaf}$. For the third, before
$\badkey \vee \badleaf$ every successful comparison lies in a
verification-first class: a class first generated by a user-$d$
encryption computation admits none, since by the final-root determinacy
of \Cref{lem:transcript-inj} and Leaf-Transcript Recovery
(\Cref{lem:leaf-recovery})---whose hypothesis $\neg\badleaf$ supplies
because the comparand is the class's encryption computation---any
non-replay submission in the class carries the same $(U,A,C)$ as that
encryption and either matches the recorded tuple, in which case the
replay check answers it before any comparison, or fails the comparison.
\Cref{lem:ver-first-root-guess} then gives
$\Pr[\forge \text{ before } \badkey \vee \badleaf]
\leq \qv^{(d)}/2^{\tauroot}$.

Summing the three bounds gives
\[
  \bigl|\Pr[\mathcal{B}^{H_{d-1}^\star} \Rightarrow 1]
       - \Pr[\mathcal{B}^{H_d^\star} \Rightarrow 1]\bigr|
  \leq
  \KeyGuess_k(\qRO(\mathcal{B}))
  + \frac{\nleaf^{(d)}}{2^{\tauleaf}}
  + \frac{\qv^{(d)}}{2^{\tauroot}},
\]
with the user-$d$ key-test event $\badkey$ charged a single time for
the adjacent hybrid.
\end{proof}
